%!TEX root = ../notes.tex
\section{April 8, 2026}
\label{20260408}

In this lecture, we will discuss private set intersection (PSI) and give a brief introduction to Fully Homomorphic Encryption. Then we give a concrete construction of Somewhat Homomorphic Encryption over Integers.

\subsection{Private Set Intersection (PSI)}

Alice and Bob want to compute the intersection of their sets $X = \{x_1, x_2, \dots, x_n\}$ and $Y = \{y_1, y_2, \cdots, y_n\}$ without revealing the elements in their set.

\begin{itemize}
    \item PSI: $f(X, Y) = X\cap Y$.
    \item PSI-Cardinality: $f(X, Y) = |X\cap Y|$ which counts the number of items in the intersection without revealing the items.
\end{itemize}

\textbf{Applications:}
\begin{itemize}
    \item \textbf{Password Breach Alert} (Chrome, Edge, Firefox, iOS Keychain, ...). There is a set of all compromised passwords, and a set of user's passwords, and we want to see if there is intersection without revealing the passwords.
    \item \textbf{Ads Conversion Measurement} (Google). Alice is an ad platform and Bob is an advertiser who has the information of users that have a made a purchase with him. The intersection is the ad conversion, those users who have seen the ads and made a purchase, which gives an idea of the effectiveness of the ad.
    \item \textbf{Privacy-Preserving Inventory Matching} (J.P. Morgan). There is a set of stocks we can sell, and a set of stocks that a buyer wants to buy. We want to which stocks we can sell to the buyer.
    \item \textbf{Private Database Joins.} For example, two companies may want to find out how many users they have in common.
\end{itemize}

\subsubsection{Na\"ive Solution}
Here's a na\"ive solution: Alice and Bob hash all their inputs, exchange hash values, and see which ones they have in common. You can learn the intersection, but could you learn \emph{more} than that?

\pseudocodeblock{
    \textbf{Alice} \< \< \textbf{Bob}\\
    \text{Input: } X = \{x_1, \dots, x_n\} \< \< \text{Input: }Y = \{y_1, \dots, y_n\}\\
    \< \sendmessageright*{H(x_1), \dots, H(x_n)} \< \\
    \< \< \text{Compute intersection between} \\
    \< \< H(x_1), \dots, H(x_n) \text{ and}\\
    \< \< H(y_1), \dots, H(y_n)
}

If the input space is a relatively small space, we can do a dictionary attack (for example, with names). This is not even semi-honest secure.

\emph{Can we even achieve 2PC/MPC with just a single round of communication (as we have done here, sending hashes one way)?} Taking a step back, we receive a message from Alice and can derive the solution regardless of \emph{any} $y$ we have. This allows us to just test multiple inputs on the function received and we'll receive that output, so we can derive $x$ from that input. Thus, no, we cannot achieve 2PC/MPC with just a single round.

\subsubsection{DDH-Based PSI}\label{sec:apr04-psi-ddh}
We start with a cyclic group $\GG$ of order $q$ with generator $g$, where DDH holds. We also have a hash function $H: \left\{ 0, 1 \right\}^{*}\to \GG$ (modeled as a random oracle).

\pseudocodeblock{
    \textbf{Alice:} \< \< \textbf{Bob:}\\
    \text{Input: } X = \{x_1, \dots, x_n\} \< \< \text{Input: }Y = \{y_1, \dots, y_n\} \\
    a \sample \ZZ_q \< \< b \sample \ZZ_q \\
    \< \sendmessageleft*{H(Y)^b := \{H(y_1)^b, \dots, H(y_n)^b\}} \< \\
    \< \sendmessageright*{H(X)^a, H(Y)^{a\cdot b}} \< \\
    \< \< \text{Compute the intersection}\\
    \< \< H(X)^{a\cdot b} \cap H(Y)^{a \cdot b}\\
}

Bob and Alice generate private exponents $k_B\sampledfrom \ZZ_q, k_A\sampledfrom \ZZ_q$ respectively. Bob will send
\[H(Y)^{k_B} := \left\{ H(y_1)^{k_B}, \cdots, H(y_n)^{k_B} \right\}\]
Alice does the same for $X$, $H(X)^{k_A}$, and sends $H(Y)^{k_A\cdots k_B}$. Bob then raises $H(X)^{k_A}$ to $k_B$ and compares. In the semi-honest case, for Bob to perform the same dictionary attack, Bob will need to break DDH in order to raise arbitrary elements $y'$ to $k_A\cdot k_B$.

We can modify this to count cardinality by randomly shuffling the returned $H(Y)^{k_A\cdot k_B}$ such that Bob cannot relate the re-encrypted hashes to the previous order.

% \Graphic{images/2023-04-04/ddh_psi-ca.png}{0.7}


\subsection{Fully Homomorphic Encryption (FHE)}

So far, our encryption schemes primarily follow an ``all-or-nothing'' idea, where if we have the secret key we can decrypt, otherwise we cannot decrypt.

In Homomorphic schemes, we want to have the additional property that an encryption of an input $x$ can be evaluated with a function $f$ to get an encryption of the output $f(x)$ without having to decrypt first.

\begin{center}
\begin{bbrenv}{fhe-box-diagram}
\begin{bbrbox}[name=Eval, namepos=center]
\end{bbrbox}
\bbrmsgto{side=Enc($x$)}
\bbrmsgto{side=$f$}
\bbrqryto{side=Enc($f(x)$)}
\end{bbrenv}
\end{center}

An additively homomorphic scheme can combine $\Enc(m_1)$ and $\Enc(m_2)$ to get
\[\Enc(m_1 + m_2)\]
as we saw in Exponential ElGamal or Paillier.

A multiplicatively homomorphic scheme can combine $\Enc(m_1)$ and $\Enc(m_2)$ to get
\[\Enc(m_1 \cdot m_2)\]
as we saw in ElGamal or RSA.

Fully homomorphic encryption means we can get both $\Enc(m_1 \cdot m_2)$ and $\Enc(m_1 + m_2)$.

\subsection{Applications}
\begin{example}[Oursourcing Storage \& Computation]
    Let's say a client stores some data on a server. A client has data $x$ and key $sk$. $ct\leftarrow\Enc(x)$ is sent to the server. If the client wants to run some computation on the server, the client sends $f$ and the server evaluates $ct' \leftarrow \mathsf{Eval}(f, ct)$ and sends $ct'$ to the client, which gives the client $f(x)$ without the server knowing $x$.
    
    \pseudocodeblock{
        \textbf{Server} \< \< \textbf{Client}\\
        \< \< \text{Data } x\\
        \< \< \text{Key sk} \\
        \< \< \text{ct} \gets \text{Enc}(x)\\
        \< \sendmessageleft*{\text{ct}, f} \< \\
        \text{ct}' \gets \text{Eval}(f, \text{ct}) \< \< \\
        \< \sendmessageright*{\text{ct}'} \< \\
        \< \< f(x) \gets \text{Dec}_{\text{sk}}(\text{ct}')
    }
\end{example}

\subsection{FHE Definition}

\begin{definition}[Homomorphic Encryption]\label{def:fhe}
    A (public-key) homomorphic encryption scheme is some
    \[\pi = (\mathsf{KeyGen}, \Enc, \Dec, \mathsf{Eval})\]
    with respect to some function family $\mathcal{F}$ with
    \begin{itemize}
        \item $(pk, sk)\leftarrow \mathsf{KeyGen}(1^\lambda)$.
        \item $ct\leftarrow \Enc_{pk}(m)\qquad m\in \{0, 1\}$.
        \item $m\leftarrow \Dec_{sk}(ct)$.
        \item $ct_f\leftarrow \mathsf{Eval}(f, ct_1, \dots, ct_n)$ with $f : \{0, 1\}^n\to \{0, 1\}$ in family $\mathcal{F}$.
    \end{itemize}

    \textbf{Correctness} requires that $\forall f\in \mathcal{F}$, if $ct_f \leftarrow \mathsf{Eval}(f, ct_1, \dots, ct_n)$ for $ct_i\leftarrow \Enc_{pk}(m_i)$, then $\Dec_{sk}(ct_f) = f(m_1, \dots, m_n)$. That is, that evaluating functions does indeed give the ciphertext of the function evaluated on the plaintexts.

    \textbf{CPA security}, as we've seen before, requires that
    \[(pk, \Enc_{pk}(m_0))\overset{c}{\simeq}(pk, \Enc_{pk}(m_1)).\]
    
    \textbf{Compactness}, that $|ct_f| \leq \mathsf{poly}(\lambda)$, that each ciphertext is bounded by some constant that is polynomial in $\lambda$ and fixed for security parameter $\lambda$. This allows us to prevent the `loophole' of evaluations returning itself.
\end{definition}

\textit{Why do we need compactness?} If we just have correctness and CPA security, our current definition is nearly trivial. For example, if our evaluation just output the tuple $(f, ct_1, \dots, ct_n)$. To decrypt, we decrypt each $ct_i$ individually and apply $f$ on top, and this satisfies our definitions! We need to add some notion that our ciphertext must stay within some size, and that we're actually \emph{combining} our ciphertexts.

If the family of functions $\mathcal{F}$ contains all polynomial sized Boolean circuits, we say that the scheme is \textbf{fully} homomorphic.

\subsubsection{Constructions}
All FHE constructions follow two steps:
\begin{enumerate}
    \item Somewhat Homomorphic Encryption (SWHE). We'll see versions over integers, from LWE (GSW), and from RLWE (BFV).
    \item Then, we bootstrap our SWHE schemes to be fully homomorphic.
\end{enumerate}

\subsubsection{SWHE over Integers}\label{sec:apr06-swhe-integers}
\textbf{Attempt 1.} This is our first attempt is a secret-key scheme:

Our secret key will be odd $p$.

\begin{remark}
    \textit{Why must $p$ be odd?} If $p$ is even, then the ciphertext $p\cdot q + m$ has the same parity as $m$. The message is not secure because we can find it just by checking even/odd. If $p$ is odd, then $p\cdot q$ can be either even or odd, so the parity of the ciphertext is not solely determined by parity of $m$.
\end{remark}

\begin{itemize}
    \item To $\Enc(m)$ with $m\in \{0, 1\}$, we sample a random $q$ and compute $ct = p\cdot q + m$. Encryption of $0$ is a multiple of $p$.
    \item Decryption is $ct\mod p$.
    \item Add: $\text{ct} \gets \text{ct}_1 + \text{ct}_2$.
    \item Multiply: $\text{ct} \gets \text{ct}_1 \cdot \text{ct}_2 $.
\end{itemize}

\emph{Is this CPA secure?} No! If we have a lot of encryptions of $0$s, taking the gcd of them will give us $p$ exactly.

\textbf{Attempt 2.} Our next attempt is to add a small error noise.

\begin{itemize}
    \item To encrypt, instead of $ct = p\cdot q + m$, we'll do
          \[ct = p\cdot q + m + 2e\]
          for small even $2e$ with $e \ll p$. Encryption of $0$ is small and even modulo $p$.
    \item Then, decryption becomes $\Dec(ct) = \left[ ct\mod p \right]\mod 2$.
    \item Addition and multiplication work the same way. Check that the decryption gives the correct result by modding p then modding 2. If
    \begin{align*}
        \text{ct}_1 &= p \cdot q_1 + m_1 + 2e_1 \\
        \text{ct}_2 &= p \cdot q_2 + m_2 + 2e_2 \\
    \end{align*}
    then
    \begin{align*}
        \text{ct}_1 + \text{ct}_2 &= p(q_1 + q_2) + (m_1 + m_2) + (2e_1 + 2e_2)\\
        \text{ct}_1 \cdot \text{ct}_2 &= pq_1(pq_2 + m_2 + 2e_2) + m_1 \cdot pq_2 + m_1 \cdot m_2\\
        & \quad + m_1 \cdot 2e_2 + 2e_1(pq_2 + m_2 + 2e_2)\\
    \end{align*}
\end{itemize}

\textbf{Attempt 3.} We can also construct a public key version of this, where the secret key is still our odd $p$ but the public key are a set of encryptions of $0$, $\{x_i = p\cdot q_i + 2e_i\}$.

Since our scheme is homomorphic over addition, we can take a random subset sum of $\{x_i\}$ and add $m$ and $2e$:
\begin{itemize}
    \item To encrypt, we'll do
          \[ct = (\text{subset sum of $\{x_i\}$}) + m + 2e\]
          for small even $2e$ with $e \ll p$. Encryption of $0$ is small and even modulo $p$.
    \item Then, decryption becomes $\Dec(ct) = \left[ ct\mod p \right]\mod 2$.
    \item Addition and multiplication work the same way.
\end{itemize}

\emph{What could go wrong?} Note that this is still only a \emph{somewhat} homomorphic encryption scheme. Are there limits to how much we can add or multiply? Specifically, could the noise grow to an extent that interferes with our encryption?
\begin{enumerate}
    \item If we add ciphertexts, our noise grows linearly:
          \begin{align*}
              ct_1        & = p\cdot q_1 + m_1 + 2e_1                         \\
              ct_2        & = p\cdot q_2 + m_2 + 2e_2                         \\
              ct_1 + ct_2 & = p\cdot (q_1 + q_2) + (m_1 + m_2) + 2(e_1 + e_2)
          \end{align*}
    \item If we multiply ciphertexts, our noise grows exponentially:
          \[ct_1 \cdot ct_2 = p\cdot (\cdots) + (m_1 \cdot m_2) + m_1\cdot 2e_2 + 2e_1m_2 + 4e_1e_2\]
\end{enumerate}
Addition is cheaper, but multiplication has our noise grow much much faster. In our parameters, we can support roughly $O(\lambda)$ multiplications, but addition is almost for free (we can do exponentially many additions).

This is why this is \emph{somewhat} homomorphic encryption---if the noise exceeds $p$, then it might affect the plaintext.

This analysis works similarly for public-key encryption\footnote{In fact, we took a very generic way of converting the secret-key scheme into a public-key scheme. We can always take subset sums of encryptions of $0$ and add it to our message.}.

Note that this protocol is quite expensive---$p$ and $q$ will tend to get quite large. We'll pivot into lattice-based cryptography which will solve some of these issues.
